\chapter{Pliego de condiciones}
\label{ch:pliego}
\section{Objetivo}
\label{sec:pliego-objetivo}
Este Pliego de Condiciones determina los requisitos a que se debe  ajustar la ejecución de las instalaciones para la distribución de energía eléctrica, cuyas características técnicas estarán especificadas en el presente pliego y correspondiente proyecto.
\section{Disposiciones generales}
\label{sec:pliego-disposiciones}
La obra deberá ajustarse a la descripción realizada en la Memoria, Planos y Presupuesto del presente proyecto.

Las calidades de los materiales deberán respetar las especificaciones mínimas.

El director técnico de la obra será la única persona capacitada para juzgar, en caso de duda y omisiones del proyecto. Lo mismo que en caso de variación de parte o del total de la obra, si no estuviese bien realizada.

El contratista está obligado al cumplimiento de la reglamentación del trabajo correspondiente, la contratación del seguro obligatorio, subsidio familiar y de vejez, seguro de enfermedad y todas aquellas reglamentaciones de carácter social vigentes o que en lo sucesivo se dicten.  

En particular deberá cumplir lo dispuesto en la norma UNE 24042 Contratación de Obras, Condiciones Generales, siempre que no modifiquen el presente Pliego de Condiciones.  

El contratista deberá estar clasificado, según Orden del Ministerio de Hacienda de 28 de Marzo de 1968 en el grupo, subgrupo y categoría correspondientes al proyecto y que se fijará en el Pliego de Condiciones Particulares, en caso de que proceda.
\section{Organización del trabajo}
\label{sec:pliego-organizacion}
El Contratista ordenará los trabajos en la forma más eficaz para la perfecta ejecución de los mismos y las obras se realizarán siempre siguiendo las indicaciones del Director de Obra, al amparo de las condiciones siguientes:
\subsection{Datos de la obra}
\label{subsec:pliego-datosOrg}
Se entregará al Contratista una copia de los planos y Pliego de Condiciones del Proyecto, así como cuantos planos o datos necesite para la completa ejecución de la Obra.

El Contratista podrá tomar nota o sacar copia a su costa de la Memoria, Presupuesto y Anexos del Proyecto, así como segundas copias de todos los documentos.

El Contratista se hace responsable de la buena conservación de los originales de donde  obtenga las copias, los cuales serán devueltos al Director de Obra después de su utilización.  

Por otra parte, en un plazo máximo de dos meses, después de la terminación de los trabajos, el Contratista deberá actualizar los diversos planos y documentos existentes, de acuerdo con las características de la Obra terminada, entregando al Director de Obra dos expedientes  completos relativos a los trabajos realmente ejecutados.  

No se harán por el Contratista alteraciones, correciones, ni adiciones o variaciones sustanciales en los datos fijados en el Proyecto, salvo aprobación previa por escrito del Director de Obra.
\subsection{Replanteo de la obra}
\label{subsec:pliego-replanteoOrg}
El Director de Obra, una vez que el Contratista esté en posesión del Proyecto y antes de comenzar las obras, deberá hacer el replanteo de las mismas, con especial atención a los puntos singulares, entregando al Contratista las referencias y datos necesarios para fijar completamente  la ubicación de las mismas.  

Se levantará por duplicado un Acta, en la que constarán, muy bien los datos entregados,  firmados por el Director de Obra y por el representante del Contratista.  

Los gastos de replanteo serán por cuenta del Contratista.   
\subsection{Mejoras y variaciones del proyecto}
\label{subsec:pliego-mejorasOrg}
No se considerarán como mejoras ni variaciones del Proyecto más que aquellas que hayan sido ordenadas expresamente por escrito,por el Director de Obra y convenido precio antes de proceder a su ejecución.  

Las obras accesorias o delicadas,  no incluidas en los precios de adjudicación,  podrán ejecutarse con personal independiente del Contratista.
\subsection{Recepción del material}
\label{subsec:pliego-recepcionOrg}
El Director de Obra, de acuerdo con el Contratista, dará a su debido tiempo su aprobación sobre el material suministrado y confirmará que permite una instalación correcta.

La vigilancia y conservación del material suministrado será por cuenta del Contratista. 
\subsection{Organización}
\label{subsec:pliego-organizacionnOrg}
El Contratista actuará de patrono legal, aceptando todas las responsabilidades correspondientes y quedando obligado al pago de los salarios y cargas que legalmente están establecidas, y en general, a todo cuanto se legisle, decrete  u ordene sobre el particular antes o durante la ejecución de la obra.

Dentro de lo estipulado en el Pliego de Condiciones, la organización de la Obra, así como la determinación de la procedencia de los materiales que se empleen, estará a cargo del Contratista a quien corresponderá la responsabilidad de la seguridad contra accidentes.  

El Contratista deberá, sin embargo, informar al Director de Obra de todos los planes de organización técnica de la Obra, así como de la procedencia de los materiales y cumplimentar  cuantas ordenes le dé éste en relación con datos extremos.  
En las obras por administración, el Contratista deberá dar cuenta diaria al Director de Obra de la admisión de personal, compra de materiales, adquisición o alquiler de elementos auxiliares y cuantos gastos haya de efectuar. 

Para los contratos de trabajo, compra de material o alquiler de elementos auxiliares, cuyos salarios, precios o cuotas sobrepasen en más de un \SI{5}{\percent} de los normales en el mercado, solicitará la aprobación previa del Director de Obra, quien deberá responder dentro de los ocho días siguientes a la petición, salvo casos de reconocida urgencia, en los que se dará cuenta posteriormente. 

\subsection{Ejecución de las obras}
\label{subsec:pliego-ejecucionOrg}
Las obras se ejecutarán conforme al Proyecto y a las condiciones contenidas en éste Pliego de condiciones y en el Pliego Particular si lo hubiera, y de acuerdo con las especificaciones señaladas en el de Condiciones Técnicas.  

El Contratista, salvo aprobación por escrito del Director de Obra, no podrá hacer ninguna alteración o modificación de cualquier naturaleza tanto en la ejecución de la obra en relación con el Proyecto, como en las Condiciones Técnicas especificadas.

El Contratista no podrá utilizar, en los trabajos, personal que no sea de su exclusiva cuenta y cargo.

Igualmente será de su exclusiva cuenta y cargo aquel personal ajeno al propiamente manual y que sea necesario para el control administratívo del mismo.

El Contratista deberá tener al frente de los trabajos un técnico suficientemente especializado a juicio del Director de Obra.
\subsection{Subcontratación de las obras}
\label{subsec:pliego-subcontrataOrg}
Salvo que el contrato disponga lo contrario o que de su naturaleza y condiciones se deduzca que la Obra ha de ser ejecutada directamente por el adjudicatario, podrá éste concertar con terceros la realización de determinadas unidades de obra.  

La celebración de los subcontratos estará sometida al cumplimiento de los siguientes requisitos:
\begin{enumerate}[label=\alph*)]
  \item A que se de conocimiento por escrito al Director de Obra del subcontrato a celebrar, con indicación de las partes de obra a realizar y sus condiciones económicas, a fin de que aquel lo autorice previamente.  
  \item A que las unidades de obra que el adjudicatario contrate con terceros no exceda del \SI{50}{\percent} del presupuesto total de la obra principal.
\end{enumerate}
En cualquier caso el Contratante no quedará vinculado en absoluto ni reconocerá ninguna obligación contractual entre él y el subcontratista y cualquier subcontratación de obras no eximirá al Contratista de ninguna de sus obligaciones con respecto al Contratante.  
\subsection{Plazo de ejecución}
\label{subsec:pliego-plazoejecucionOrg}
Los plazos de ejecucion, total y parciales, indicados en el contrato, se empezarán a contar a partir de la fecha de replanteo.  
El Contratista estará obligado a cumplir con los plazos que se señalen en el contrato para la ejecución de las obras y que serán improrrogables.   

No obstante lo anteriormente indicado, los plazos podrán ser objeto de modificaciones cuando así resulte por cambios determinados por el Director de Obra debidos a exigencias de la realización de las obras y siempre que tales cambios influyan realmente en los plazos señalados en el contrato.   

Si por cualquier causa, ajena por completo al Contratista, no fuera posible empezar los trabajos en la fecha prevista o tuvieran que ser suspendidos una  vez empezados, se concederá por el Director de Obra, la prorroga estrictamente necesaria.  
\subsection{Recepción provisional}
\label{subsec:pliego-recepcionprovisionalOrg}
Una vez terminadas las obras y a los quince días siguientes a la petición del Contratista se hará la recepción provisional de las mismas por el Contratante, requiriendo para ello la presencia del Director de Obra y del representante del Contratista levantandose las Actas que correspondan en las que se harán constar la conformidad con los trabajos realizados, si éste es el caso.

Dichas Actas serán firmadas por el Director de Obra y el representante del Contratista, dándose la Obra por recibida si se ha ejecutado correctamente de acuerdo con las especificaciones dadas en el Pliego de Condiciones Técnicas y en el Proyecto correspondiente, comenzandose entonces a contar el plazo de garantía.

En el caso de no hallarse la Obra en estado de ser recibidad,  se hará constar  así en el Acta y  se darán al Contratista las instrucciones precisas y detalladas para remediar los defectos observados, fijandose un plazo de ejecución.

Expirado dicho plazo, se hará un nuevo reconocimiento.  Las obras de reparación  serán por cuenta y a cargo del Contratista.  
Si el Contratista no cumpliese estas prescripciones podrá declararse rescindido  el contrato con pérdida de la fianza.  
\subsection{Períodos de garantía}
\label{subsec:pliego-garantiaOrg}
El periodo de garantía será señalado en el contrato y empezará a contar desde la fecha de aprobación del Acta de Recepción. 

Hasta que tenga lugar la recepción definitiva, el Contratista es responsable de la conservación de la Obra, siendo de su cuenta y cargo las reparaciones por defectos de ejecución o mala calidad de los materiales.

Durante  este periodo, el Contratista garantizará al Contratante contra toda  reclamación de terceros, fundada en causa y por ocasión de la ejecución de la Obra.
\subsection{Recepción definitiva}
\label{subsec:pliego-recepciondefOrg}
Al terminar el Plazo de garantía señalado en el contrato o en su defecto a los  seis meses de la recepción provisional, se procedera a la recepción definitiva  de las obras, con la concurrencia del Director de Obra y del representante del Contratista levantándose el Acta correspondiente, por duplicado (si las obras son conformes), que quedará firmada por el Director de Obra y el representante  del Contratista y ratificada por el Contratante y el Contratista.
\subsection{Pago de obras}
\label{subsec:pliego-pagoOrg}
El pago de las obras realizadas se hará sobre certificaciones parciales, que se practicarán mensualmente. Dichas certificaciones contendrán solamente las unidades de obra totalmente terminadas que se hubieran ejecutado en el plazo a que se refieran.  

La relación valorada que figure en las certificaciones, se hará con arreglo a los precios establecidos, y con la ubicación, planos y referencias necesarias para su comprobación.

El Director de Obra expedirá las Certificaciones de las obras ejecutadas que tendrán carácter de documento provisional a buena cuenta, rectificables por la  liquidación definitiva o por las certificaciones siguientes.
\subsection{Abono de materiales acopiados}
\label{subsec:pliego-abonomatOrg}
Cuando a juicio del Director de Obra no haya peligro de que desaparezcan o se  deterioren los materiales acopiados y reconocidos como útiles, se abonarán con arreglo a los precios descompuestos de la adjudicación. 

Dicho material será indicado por el Director de Obra e indicado en el Acta de recepción de Obra.  

La restitución de las bobinas vacías se hará en el plazo de un mes, una vez que  se haya instalado el cable que contenían.  
\section{Condiciones técnicas de la ejecución}
\label{sec:pliego-condiciones}
\subsection{Excavaciones}
\label{subsec:pliego-excavacionesCT}
Las dimensiones de las excavaciones se ajustarán lo más posible a las dadas en el Proyecto o en su defecto a las indicadas por el Director de Obra.  

Las paredes de los hoyos serán verticales. Cuando sea necesario variar el volumen de la excavación, se hará de acuerdo con el Director de Obra.  

El Contratista tomara las disposiciones convenientes para dejar el menor tiempo posible abiertas las excavaciones, con objeto de evitar accidentes. Las excavaciones se realizarán con útiles apropiados según el tipo de terreno.   

En terrenos rocosos será imprescindible el uso de explosivos o martillo compresor, siendo por cuenta del Contratista la obtención de los permisos de  utilización de explosivos.  

Cuando deban emplearse explosivos, el Contratista deberá tomar las precauciones adecuadas para que en el momento de la explosión no se proyecten al exterior piedras que puedan provocar accidentes o desperfectos, cuya responsabilidad correría a cargo del Contratista.

En terrenos con agua deberá procederse a su desecado, procurando hormigonar después lo más rapidamente posible para evitar el riesgo de desprendimientos en las paredes del hoyo, aumentando así las dimensiones del mismo.

\subsection{Hormigonado}
\label{subsec:pliego-hormigonadoCT}
Este se deberá dosificar a \SI{250}{\kilo\gram} de cemento por cada metro cúbico.

Si la excavación superara el \SI{10}{\percent} del volumen técnico, por conveniencia del contratista, siempre de acuerdo con el Director técnico de las obras, o el empleo de explosivos, la dosificación del hormigón será siempre la misma.

El cemento empleado será Portland, de fraguado lento, o bien de otra marca similar, de primera calidad.
    
Los áridos empleados para las cimentaciones de los apoyos, deberán ser de buena calidad, limpios y no heladizos, estando exentos de materiales orgánicos y de arcillas.

Será preferible la piedra con aristas y superficies rugosas y ásperas, por su mayor adherencia al mortero.

La arena puede proceder de minas o canteras, ríos, o bien, de machaqueo.

La dimensión de los granos de arena no será superior al \SI{6}{\percent} (ensayo de granulometría).

El agua empleada para la ejecución del hormigón será limpia y exenta de elementos orgánicos, arcillas, etc.
\subsection{Armado e izado de los apoyos metálicos}
\label{subsec:pliego-izadoCT}
El transporte de todos los materiales a la obra se realizará con el mayor cuidado, e intentando evitar al máximo los posibles desperfectos que pudieran acontecer.

En caso de dobleces de barras, éstas se enderezarán en caliente. Los taladros que se tengan que realizar, se harán con punzón o carraca, nunca por sopletes. Los taladros que no se usen, se cerrarán por medio de soldadura. En caso de que haya que aumentar el diámetro de los mismos, se hará por mediación del escariador. Se deberán eliminar las rebabas de los mismos.

Para el armado se empleará puntero y martillo para que coincidan las piezas que se unen, pero con cuidado para no agrandar el taladro.
Se aconseja armar en tierra el mayor número posible de piezas.

El izado deberá hacerse sin originar deformaciones permanentes sobre elementos que componen el apoyo.
Cuando la torre está izada, se hará un repaso general del ajuste de los componentes.

Los postes de hormigón se transportarán en vehículos preparados al efecto, y, al depositarlos se hará en un lugar llano y con sumo cuidado en evitación de deformaciones de los mismos.

Todas las piezas deberán estar recubiertas de material blando y flexible (gomas naturales o sintéticas).
\subsection{Tendido, tensado y regulado de los conductores}
\label{subsec:pliego-tendidoCT}
Los cables deberán tratarse con el mayor cuidado para evitar deterioros, lo mismo que las bobinas donde se transportan.

En la hora de desenrollar los cables se debe cuidar que no rocen con el suelo.

Para ejercer la tracción se pueden emplear cuerdas pilotos, pero deben ser las mismas del tipo flexible y antigiratorias, montando bulones de rotación para compensar los defectos de la torsión. Si se produce alguna rotura en los hilos de los cables, por cualquier causa, se deberán colocar manguitos separatorios.
	
Todo el tendido y tensado de los conductores se realizará conforme a la tabla de tendido proporcionada por el proyectista, y conforme a las características climatológicas a las que se va a realizar la operación.
\begin{itemize}
    \item \textbf{Poleas de tendido:} Para cables de aluminio, éstas serán de aleación de aluminio. El diámetro será entre 25 y 30 veces el diámetro del cable que se extienda. Esta polea estará calculada para aguantar esfuerzos a que deba ser sometida.
    \item \textbf{Tensado:} Este deberá realizarse arriostrando las torres de amarre a los apoyos de hormigón de anclajes en sentido longitudinal. El tensado de los cables se hará por medio de un cable piloto de acero en evitación de flexiones exageradas. Todos los aparatos para el tensado deberán colocarse a distancia conveniente de la torre de tense, para que el ánguo formado por las tangentes del piloto al paso por la polea no sea inferior a os 150 grados.
    \item \textbf{Regulado:} Toda línea se divide en trozos de longitudes varialbles según situación de vértices. En el perfil longitudinal se definen los vanos y en los cálculos las flechas de cada uno de ellos, y al mismo se deberá adaptar.
\end{itemize}
\subsection{Cadena de aisladores}
\label{subsec:pliego-aislCT}
Estos se limpiarán cuidadosamente antes de ser montados. Se tendrá especial cuidado en su traslado y colocación para que no sufran desperfectos los herrajes que unen las cadenas.
\subsection{Empalmes}
\label{subsec:pliego-empalmesCT}
Serán de tal calidad que garanticen la resistencia mecánica exigida por los Reglamentos y no exista aumento de la resistencia del conductor.

Los empalmes deberán ser cepillados cuidadosamente, tanto interior como exteriormente, con cepillo y baquetas especiales.
\subsection{Engrapado}
\label{subsec:pliego-engrapadoCT}
Para el mismo se deberá tomar medida para conseguir un buen aplomo de las cadenas de aisladores.
	
El apretado de los tornillos de las grapas se debe hacer alternativamente para asegurar un buen apriete.
\section{Características de los materiales}
\label{sec:pliego-materiales}
Todos los materiales serán de primera calidad. No deberán presentar deterioro ni defecto alguno que disminuya la función que tengan que desarrollar.
\subsection{Conductores trenzados}
\label{subsec:pliego-trenzadosMAT}
Deberán ir provistos de cubierta de aislamiento, el cual será de polietileno reticulado (PRC).

Se deberán distinguir de otros por lo que deberán ir grabados en tintas blancas o relieves en el exterior.

Las secciones de los conductores serán las determinadas en la Memoria.

Los empalmes deberán realizarse mediante manguitos a compresión y el aislamiento será regenerado con cinta de goma autovulcanizante y recubierta con cinta de P.V.C.
\subsection{Conductores de cobre}
\label{subsec:pliego-cobreMAT}
Estos estarán formados, según la sección, por uno o por varios alambres de cobre, cilíndricos, de buena calidad y resistencia mecánica y libres de todos los desperfectos posibles, así como de imperfecciones.
\subsection{Abrazaderas y tacos de sujeción}
\label{subsec:pliego-abrazaderasMAT}
Las abrazaderas serán de placas de acero isoplastificados y de una sola pieza, dotadas de punta de acero roscada.

Las abrazaderas para cable fiador, serán las mismas, de iguales características, pero sin punta de acero.
\subsection{Herrajes}
\label{subsec:pliego-herrajesMAT}
El cable fiador de acero y de arriostramiento será flexible y galvanizado.

El resto de los herrajes (aprietahilos, grilletes, etc.), serán galvanizados en caliente.
\subsection{Torres metálicas}
\label{subsec:pliego-torresMAT}
Serán de hierro laminado y responderán a la altura determinada en la Memoria.

Serán galvanizadas en caliente. Las cimentaciones se tendrán que adaptar a lo especificado en el cálculo de las mismas.


